\section{子 Agent：Claude Code 的分身术}

\begin{knowledgebox}{子 Agent 机制}
子 Agent 是独立的 AI 进程，每个有自己的上下文窗口，不会污染主对话。内置三种类型：
\begin{itemize}[leftmargin=1.5em]
  \item \textbf{Explore}：只读，快速搜索代码库，适合找文件、找定义
  \item \textbf{Plan}：只读，研究分析，用于 Plan Mode 下的代码分析
  \item \textbf{General}：完整能力，适合复杂的多步骤任务
\end{itemize}
\end{knowledgebox}

\subsection{核心价值：隔离上下文}

当你让 Claude 跑测试、分析日志、搜索大量文件时，海量输出会塞满上下文窗口。用子 Agent 来做这些事，输出留在子 Agent 的上下文里，只有摘要返回给你。

按 \texttt{Ctrl+B} 可以把子 Agent 放到后台运行——适合跑测试套件、大范围代码搜索、不紧急的代码审查。

\subsection{自定义子 Agent}

在 \texttt{.claude/agents/} 目录下创建自定义 Agent 的 Markdown 文件，然后在对话中用 \texttt{@"agent-name (agent)"} 调用。可以为特定场景（如 code review）定制专属 Agent。

\subsection{本章小结}

子 Agent 的核心价值是上下文隔离。把高噪声操作（测试、日志分析、大范围搜索）交给子 Agent，保持主对话干净高效。

%% ====================================================================
